\section{Extension of scalars and splitting fields}

    Let \(R\) be a commutative ring, let \(A\) be an \(R\)-algebra, and let \(U\)
    be an \(A\)-module. Suppose that
    \[
        \varphi:R\longrightarrow R'
    \]
    is a homomorphism of commutative rings. Then \(R'\) is naturally an
    \(R\)-algebra via \(\varphi\).

\begin{definition}[Extension of scalars]
    With the notation above, the tensor product
    \[
        R'\otimes_R A
    \]
    is naturally an \(R'\)-algebra. Moreover,
    \[
        R'\otimes_R U
    \]
    is naturally a module over the \(R'\)-algebra \(R'\otimes_R A\).

    Explicitly, the action is given on elementary tensors by
    \[
        (r'\otimes a)\cdot (s'\otimes u)
        :=
        r's'\otimes au,
    \]
    where \(r',s'\in R'\), \(a\in A\), and \(u\in U\).
\end{definition}

\begin{example}
    Let
    \[
        R=\mathbb{Q},
        \qquad
        R'=\mathbb{C},
        \qquad
        A=M_n(\mathbb{Q}).
    \]
    Then
    \[
        \mathbb{C}\otimes_{\mathbb{Q}}M_n(\mathbb{Q})
        \cong M_n(\mathbb{C}).
    \]
    Thus extending scalars from \(\mathbb{Q}\) to \(\mathbb{C}\) turns the
    \(\mathbb{Q}\)-algebra \(M_n(\mathbb{Q})\) into the \(\mathbb{C}\)-algebra
    \(M_n(\mathbb{C})\).
\end{example}

\begin{remark}
    If \(R\subseteq R'\) is a subring and \(U\) is an \(A\)-module, then we say that
    the module
    \[
        V:=R'\otimes_R U
    \]
    is obtained from \(U\) by extending scalars from \(R\) to \(R'\).
\end{remark}

\begin{definition}
    Let \(V\) be an \(R'\otimes_R A\)-module. If there exists an \(A\)-module \(U\)
    such that
    \[
        V\cong R'\otimes_R U
    \]
    as \(R'\otimes_R A\)-modules, then we say that \(V\) can be written over \(R\),
    or that \(V\) is obtained from \(U\) by extension of scalars.
\end{definition}

\begin{example}
    Let \(R\) be a discrete valuation ring, and let \(K=\operatorname{Frac}(R)\) be
    its fraction field. If \(A\) is an \(R\)-algebra which is finitely generated as
    an \(R\)-module, then
    \[
        K\otimes_R A
    \]
    is a finite-dimensional \(K\)-algebra.
\end{example}

\begin{proof}
    Since \(A\) is finitely generated as an \(R\)-module, there exist elements
    \[
        a_1,\dots,a_m\in A
    \]
    such that
    \[
        A=Ra_1+\cdots+Ra_m.
    \]
    We first show that \(K\otimes_R A\) is finite-dimensional as a \(K\)-vector
    space.

    The \(K\)-vector space structure on \(K\otimes_R A\) is given by
    \[
        \lambda\cdot(\mu\otimes a):=(\lambda\mu)\otimes a,
        \qquad
        \lambda,\mu\in K,\ a\in A.
    \]
    Every element of \(K\otimes_R A\) is a finite sum of elementary tensors
    \[
        \lambda\otimes a,
        \qquad \lambda\in K,\ a\in A.
    \]
    Since \(a_1,\dots,a_m\) generate \(A\) as an \(R\)-module, for each \(a\in A\)
    there exist \(r_1,\dots,r_m\in R\) such that
    \[
        a=r_1a_1+\cdots+r_ma_m.
    \]
    Hence
    \[
        \lambda\otimes a
        =
        \lambda\otimes \left(\sum_{i=1}^m r_i a_i\right)
        =
        \sum_{i=1}^m \lambda r_i\otimes a_i.
    \]
    Since
    \[
        \lambda r_i\otimes a_i
        =
        (\lambda r_i)\cdot (1\otimes a_i),
    \]
    it follows that
    \[
        1\otimes a_1,\dots,1\otimes a_m
    \]
    span \(K\otimes_R A\) as a \(K\)-vector space. Therefore
    \[
        \dim_K(K\otimes_R A)\leq m<\infty.
    \]

    It remains to describe the algebra structure. Since \(A\) is an \(R\)-algebra
    and \(K\) is an \(R\)-algebra via the inclusion
    \[
        R\hookrightarrow K,
    \]
    the tensor product \(K\otimes_R A\) carries a natural multiplication defined on
    elementary tensors by
    \[
        (\lambda\otimes a)(\mu\otimes b)
        :=
        \lambda\mu\otimes ab,
        \qquad
        \lambda,\mu\in K,\ a,b\in A.
    \]
    This multiplication is \(K\)-bilinear and associative, and its identity element is
    \[
        1_K\otimes 1_A.
    \]
    Hence \(K\otimes_R A\) is a \(K\)-algebra.

    Since it is finite-dimensional as a \(K\)-vector space, we conclude that
    \[
        K\otimes_R A
    \]
    is a finite-dimensional \(K\)-algebra.
\end{proof}

\begin{remark}
    Suppose that \(V=R'\otimes_R U\) can be written over \(R\). If \(U\) is free as
    an \(R\)-module, then we identify \(U\) with its image
    \[
        1_{R'}\otimes_R U\subseteq R'\otimes_R U.
    \]
    In this case, an \(R\)-basis of \(U\) becomes an \(R'\)-basis of
    \(R'\otimes_R U\).
\end{remark}

\begin{example}[Reduction modulo an ideal]
    Let \(I\) be an ideal of \(R\), and take
    \[
        R'=R/I.
    \]
    Then applying the functor
    \[
        R'\otimes_R -
    \]
    to an \(R\)-module \(U\) is the same as reducing \(U\) modulo \(I\). More
    precisely,
    \[
        (R/I)\otimes_R U\cong U/IU.
    \]
\end{example}

\begin{definition}[Lifting of modules]
    Let \(V\) be an \(R'\otimes_R A\)-module. If there exists an \(A\)-module \(U\)
    such that
    \[
        V\cong R'\otimes_R U,
    \]
    then we say that \(V\) can be lifted to \(U\), and that \(U\) is a lift of \(V\).
\end{definition}

\begin{example}
    Let \(R\) be a discrete valuation ring with maximal ideal
    \[
        (\pi)=J(R),
    \]
    and let \(A\) be an \(R\)-algebra. Put
    \[
        F=R/(\pi).
    \]
    Then
    \[
        F\otimes_R A\cong A/\pi A.
    \]
    If \(V\) is an \(A\)-module, then
    \[
        F\otimes_R V\cong V/\pi V.
    \]
    Thus tensoring with the residue field \(F\) corresponds to reducing modulo
    \(\pi\).
\end{example}

\begin{proposition}
    Let
    \[
        F\subseteq E
    \]
    be an algebraic extension of fields. Let \(A\) be a finite-dimensional
    \(F\)-algebra, and let \(V\) be an \(E\otimes_F A\)-module which is finite-dimensional
    as an \(E\)-vector space.

    Then there exists a field \(K\) such that
    \[
        F\subseteq K\subseteq E,
        \qquad
        [K:F]<\infty,
    \]
    and \(V\) can be written over \(K\). More precisely, there exists a
    \(K\otimes_F A\)-module \(V_K\) such that
    \[
        V\cong E\otimes_K V_K
    \]
    as \(E\otimes_F A\)-modules.
\end{proposition}

\begin{proof}
    If \(V=0\), the result is trivial. Assume \(V\neq 0\), and write
    \[
        d=\dim_E V.
    \]
    Choose an \(E\)-basis
    \[
        v_1,\dots,v_d
    \]
    of \(V\).

    Let
    \[
        a^1,\dots,a^n
    \]
    be an \(F\)-basis of \(A\). Since \(V\) is an \(E\otimes_F A\)-module, each element
    \(a^t\in A\) acts \(E\)-linearly on \(V\). With respect to the chosen \(E\)-basis of
    \(V\), write the action of \(a^t\) as
    \[
        a^t v_j=\sum_{i=1}^d a_{ij}^t v_i,
        \qquad
        a_{ij}^t\in E.
    \]
    Define
    \[
        K:=F\bigl(a_{ij}^t\mid 1\leq t\leq n,\ 1\leq i,j\leq d\bigr)
        \subseteq E.
    \]
    Since \(E/F\) is algebraic, each \(a_{ij}^t\) is algebraic over \(F\). Hence the
    field obtained by adjoining finitely many such elements is a finite extension of
    \(F\). Therefore
    \[
        [K:F]<\infty.
    \]

    Now define
    \[
        V_K:=Kv_1+\cdots+Kv_d.
    \]
    This is a \(K\)-vector space with basis \(v_1,\dots,v_d\). By construction, for
    every \(t\) and every \(j\),
    \[
        a^t v_j=\sum_{i=1}^d a_{ij}^t v_i\in V_K,
    \]
    because all coefficients \(a_{ij}^t\) lie in \(K\). Thus \(V_K\) is stable under
    the action of \(A\).

    Hence \(V_K\) becomes a \(K\otimes_F A\)-module by
    \[
        (\lambda\otimes a)\cdot v
        :=
        \lambda(av),
        \qquad
        \lambda\in K,\ a\in A,\ v\in V_K.
    \]

    Finally, consider the natural \(E\)-linear map
    \[
        E\otimes_K V_K\longrightarrow V,
        \qquad
        \lambda\otimes v\longmapsto \lambda v.
    \]
    Since \(v_1,\dots,v_d\) is both a \(K\)-basis of \(V_K\) and an \(E\)-basis of
    \(V\), this map is an isomorphism of \(E\)-vector spaces.

    Moreover, it is compatible with the \(E\otimes_F A\)-module structures. Indeed,
    for \(\mu,\lambda\in E\), \(a\in A\), and \(v\in V_K\),
    \[
        (\mu\otimes a)\cdot(\lambda\otimes v)
        =
        \mu\lambda\otimes av
        \longmapsto
        \mu\lambda av,
    \]
    which is exactly the action of \(\mu\otimes a\) on \(\lambda v\in V\).

    Therefore
    \[
        V\cong E\otimes_K V_K
    \]
    as \(E\otimes_F A\)-modules. Hence \(V\) can be written over the finite extension
    \(K/F\).
\end{proof}

\begin{definition}[Absolutely simple module]
    Let \(F\) be a field, and let \(A\) be a finite-dimensional \(F\)-algebra.
    A simple \(A\)-module \(U\) is called absolutely simple if, for every field
    extension \(E/F\), the \(E\otimes_F A\)-module
    \[
        E\otimes_F U
    \]
    is simple.
\end{definition}

\begin{definition}[Splitting field]
    Let \(F\) be a field, and \(A\) be a finite-dimensional \(F\)-algebra.
    A field extension \(E/F\) is called a splitting field for \(A\) if every simple
    \(E\otimes_F A\)-module is absolutely simple.
\end{definition}

\begin{proposition}
    Let \(A\) be a finite-dimensional algebra over a field \(F\), and let \(U\) be
    a simple \(A\)-module. Then the following conditions are equivalent:
    \begin{enumerate}
        \item \(U\) is absolutely simple.
        \item
        \[
            \operatorname{End}_A(U)=F.
        \]
        \item The matrix algebra summand of \(A/J(A)\) corresponding to \(U\) has
        the form
        \[
            M_n(F),
        \]
        where
        \[
            n=\dim_F U.
        \]
    \end{enumerate}
\end{proposition}

\begin{proof}
    We prove the equivalences.

    \((1)\Rightarrow (2)\).
    Let \(\overline{F}\) be an algebraic closure of \(F\). Since \(U\) is absolutely
    simple,
    \[
        \overline{F}\otimes_F U
    \]
    is a simple \(\overline{F}\otimes_F A\)-module.

    By Schur's lemma,
    \[
        \operatorname{End}_{\overline{F}\otimes_F A}
        \bigl(\overline{F}\otimes_F U\bigr)
    \]
    is a division algebra over \(\overline{F}\). Since \(\overline{F}\) is algebraically
    closed and \(\overline{F}\otimes_F U\) is finite-dimensional over \(\overline{F}\),
    we get
    \[
        \operatorname{End}_{\overline{F}\otimes_F A}
        \bigl(\overline{F}\otimes_F U\bigr)
        =
        \overline{F}.
    \]
    On the other hand, by base change for endomorphism rings,
    \[
        \operatorname{End}_{\overline{F}\otimes_F A}
        \bigl(\overline{F}\otimes_F U\bigr)
        \cong
        \overline{F}\otimes_F \operatorname{End}_A(U).
    \]
    Therefore
    \[
        \overline{F}\otimes_F \operatorname{End}_A(U)
        \cong
        \overline{F}.
    \]
    Taking dimensions over \(\overline{F}\), we obtain
    \[
        \dim_F\operatorname{End}_A(U)=1.
    \]
    Since \(F\) embeds into \(\operatorname{End}_A(U)\) by scalar multiplication, it
    follows that
    \[
        \operatorname{End}_A(U)=F.
    \]

    \((2)\Rightarrow (3)\).
    By the Artin--Wedderburn theorem, the simple module \(U\) corresponds to a matrix
    algebra summand of \(A/J(A)\) of the form
    \[
        M_n(D),
    \]
    where
    \[
        D=\operatorname{End}_A(U)^{\operatorname{op}}.
    \]
    If
    \[
        \operatorname{End}_A(U)=F,
    \]
    then
    \[
        D=F.
    \]
    Hence the corresponding summand is
    \[
        M_n(F).
    \]
    In this case the natural simple module for \(M_n(F)\) is \(F^n\), and therefore
    \[
        n=\dim_F U.
    \]

    \((3)\Rightarrow (1)\).
    Suppose that the matrix algebra summand of \(A/J(A)\) corresponding to \(U\) is
    \[
        M_n(F),
    \]
    with
    \[
        n=\dim_F U.
    \]
    Then \(A\) acts on \(U\) through a surjective algebra homomorphism
    \[
        A\longrightarrow M_n(F),
    \]
    and under this action we may identify
    \[
        U\cong F^n.
    \]

    Let \(E/F\) be any field extension. After extending scalars, \(E\otimes_F A\) acts
    on \(E\otimes_F U\) through the induced homomorphism
    \[
        E\otimes_F A
        \longrightarrow
        E\otimes_F M_n(F)
        \cong
        M_n(E).
    \]
    Moreover,
    \[
        E\otimes_F U
        \cong
        E\otimes_F F^n
        \cong
        E^n.
    \]
    The natural \(M_n(E)\)-module \(E^n\) is simple. Therefore
    \[
        E\otimes_F U
    \]
    is a simple \(E\otimes_F A\)-module for every field extension \(E/F\).

    Hence \(U\) is absolutely simple.
\end{proof}

\begin{example}
    Let
\[
    f(x)=x^4-2\in \mathbb{Q}[x],
\]
and let
\[
    \alpha=\sqrt[4]{2},\qquad i=\sqrt{-1}.
\]
The roots of \(f(x)\) in \(\mathbb{C}\) are
\[
    \alpha,\quad i\alpha,\quad -\alpha,\quad -i\alpha.
\]
Thus the splitting field of \(x^4-2\) over \(\mathbb{Q}\) is
\[
    K=\mathbb{Q}(\alpha,i).
\]
Since \(x^4-2\) is irreducible over \(\mathbb{Q}\) and \(i\notin \mathbb{Q}(\alpha)\), we have
\[
    [K:\mathbb{Q}]=8.
\]
Therefore
\[
    \left|\operatorname{Gal}(K/\mathbb{Q})\right|=8.
\]

The Galois group acts on the set of roots
\[
    \Omega=\{\alpha,i\alpha,-\alpha,-i\alpha\}.
\]
Consider the automorphisms \(\sigma,\tau\in \operatorname{Gal}(K/\mathbb{Q})\) defined by
\[
    \sigma(\alpha)=-i\alpha,\qquad \sigma(i)=-i,
\]
and
\[
    \tau(\alpha)=\alpha,\qquad \tau(i)=-i.
\]
Their actions on the roots are
\[
    \sigma:\quad
    \alpha\leftrightarrow -i\alpha,\qquad
    -\alpha\leftrightarrow i\alpha,
\]
and
\[
    \tau:\quad
    i\alpha\leftrightarrow -i\alpha,
    \qquad
    \alpha\mapsto \alpha,\qquad -\alpha\mapsto -\alpha.
\]
Hence
\[
    \sigma^2=\tau^2=(\sigma\tau)^4=1.
\]
Moreover,
\[
    \langle \sigma,\tau\rangle=\operatorname{Gal}(K/\mathbb{Q}).
\]
Therefore
\[
    \operatorname{Gal}(K/\mathbb{Q})
    \cong
    D_8,
\]
where
\[
    D_8
    =
    \left\langle \sigma,\tau
    \mid
    \sigma^2=\tau^2=(\sigma\tau)^4=1
    \right\rangle
\]
is the dihedral group of order \(8\).

The group \(D_8\) can be represented as a subgroup of
\[
    \operatorname{GL}_2(\mathbb{R}).
\]
Indeed, consider the real vector space
\[
    E=\mathbb{R}\alpha\oplus \mathbb{R}i\alpha
\]
with ordered basis
\[
    (\alpha,i\alpha).
\]
With respect to this basis, the above actions are represented by the matrices
\[
    \sigma\longmapsto
    \begin{pmatrix}
        0 & -1\\
        -1 & 0
    \end{pmatrix},
    \qquad
    \tau\longmapsto
    \begin{pmatrix}
        1 & 0\\
        0 & -1
    \end{pmatrix}.
\]
These matrices satisfy
\[
    \sigma^2=\tau^2=I_2,
    \qquad
    (\sigma\tau)^4=I_2.
\]
Thus they generate a subgroup of \(\operatorname{GL}_2(\mathbb{R})\) isomorphic to
\(D_8\).

Geometrically, this is the usual symmetry group of a square whose vertices are
\[
    \alpha,\quad i\alpha,\quad -\alpha,\quad -i\alpha.
\]
\end{example}
